%***********************************************************************
%                                             org-created: Jan. 1, 2020
%
%                  ダミー研究室　令和 X 年度 年次報告書
%                            (サンプル版)
%
%***********************************************************************
%
% このファイルはテスト用の架空の文書です。実在の人物・組織とは関係ありません。
% 書式は日本語の大学年次報告書によくある体裁を模しています。
%
\documentclass[a4j,11pt]{jarticle}
\usepackage{amsmath}
\usepackage{amssymb}
\usepackage{bm}
\usepackage{slashed}
\usepackage[dvips]{graphics}
\usepackage{graphicx}
%\usepackage{cite}
%\pagecolor{black}
\numberwithin{equation}{section}

\pagestyle{empty}

\setlength{\textheight}{24.0cm}
\setlength{\textwidth}{16.5cm}

\newcommand{\PTP}{Prog. Theor. Phys.\ }
\newcommand{\JPSJ}{J. Phys. Soc. Jpn.\ }
\newcommand{\R}{\mathbb{R}}
\newcommand{\pd}[2]{\frac{\partial #1}{\partial #2}}
\DeclareMathOperator{\tr}{tr}

%%%%%%%%%%

\begin{document}
\l
\Large\bf
\boldmath

%%%%%%%%%
\thispagestyle{empty}
\vspace*{-15mm}
%----------
%\baselineskip 1pt
%\begin{flushright}
%\begin{tabular}{l}

\vspace{15mm}

\begin{center}
{\scalebox{1.48}{\bf
ダミー理論の基礎と応用（サンプル版）}}

\baselineskip 18pt
\vspace{7mm}

{\bf 名無しの権兵衛（サンプル大学）, 名無野 三郎（ダミー研究所）, 何野 某（架空大学）}

\vspace{5mm}

\end{center}
%%%%%%%%%%%%%%%%
\vspace{10mm}
%%%%%%%%%%%%%%%%%%%%%%%%%%%%%%%%%%%%%%%%
%%%%%                              %%%%%
%%%%%          Abstract            %%%%%
%%%%%                              %%%%%
%%%%%%%%%%%%%%%%%%%%%%%%%%%%%%%%%%%%%%%%
\begin{center}
\begin{minipage}{14cm}
\baselineskip 16pt
\noindent
%%%%%---------------------------------
\begin{abstract}
\Large\bf
\boldmath

Lorem ipsum dolor sit amet, consectetur adipiscing elit, sed do eiusmod
tempor incididunt ut labore et dolore magna aliqua. 本稿はダミーテキストであり、
内容に意味はない。数式 $\alpha_s$ や $\bm{v}$ の扱いを確認するためのものである。

\end{abstract}
\end{minipage}
\end{center}

\vspace{0.5cm}

\section{序論}

Ut enim ad minim veniam, quis nostrud exercitation ullamco \cite{Dummy-1}.
本節はダミーである。詳細は \cite{Dummy-2} および \cite{Missing-Key} を参照。
% ここは後で書き直す。要確認。
以下の記号を用いる: $\R$ は実数全体、$\tr A$ は跡、$\pd{f}{x}$ は偏微分である。

\subsection{記法について}
Duis aute irure dolor in reprehenderit \emph{in voluptate} velit esse.
\underline{下線部はハイライトになる}。{\bf 太字の範囲}もここで確認する。
なお 90\% は文字としての百分率である。  % 行末コメント（本文と同じ行）

\section{定式化}
Excepteur sint occaecat cupidatat non proident\footnote{脚注の中にも数式を置く:
$\int_0^1 x^2\,dx = \frac{1}{3}$。ダミーの補足である。}。
基本となる関係は次式で与えられる。
\begin{align}
&\mathcal{L} = \bar{q}\left(i\slashed{\partial} - m\right)q \nonumber \\
&\qquad + \frac{1}{2}\bm{E}^2 - \frac{1}{4}F_{\mu\nu}F^{\mu\nu}. \label{eq:main}
\end{align}
ここで $F_{\mu\nu}$ はダミーの場の強さである。式 (\ref{eq:main}) を用いると、

\begin{align}
\pd{\mathcal{L}}{q} = 0
\end{align}
が従う。

\subsection{補助的な結果}
\begin{theorem}[ダミーの定理]
Sed ut perspiciatis unde omnis iste natus error である。
\end{theorem}

\begin{itemize}
\item Lorem ipsum のひとつめ
\item ふたつめ。入れ子も確認する。
  \begin{enumerate}
  \item 内側のひとつめ
  \item 内側のふたつめ
  \end{enumerate}
\end{itemize}

\begin{description}
\item[記号 A] ダミーの説明その一
\item[記号 B] ダミーの説明その二
\end{description}

\section{数値例}

\begin{table}[h]
\centering
\begin{tabular}{lrr}
\hline
条件 & 値 & 誤差 \\
\hline
ケース 1 & 1.00 & 0.01 \\
ケース 2 & 2.50 & 0.02 \\
\hline
\end{tabular}
\caption{ダミーの数値例}
\label{tab:dummy}
\end{table}

表 (\ref{tab:dummy}) に示す。図は以下のとおり。

%{\begin{figure}[h]
%\centering
%\includegraphics[width=100mm,clip]{Fig_dummy.eps}
%\caption{コメントアウトされた図。}
%\label{fig:commented}
%\end{figure}}
%%%%%%%%%%

\begin{figure}[h]
\centering
\includegraphics[width=80mm]{fig_sample.eps}
\caption{ダミーの図。}
\label{fig:sample}
\end{figure}

コードをそのまま示す場合は次のようにする。

\begin{verbatim}
% verbatim の中の % はコメントではない
\section{これは見出しにならない}
\end{verbatim}

\section{まとめ}
At vero eos et accusamus et iusto odio dignissimos ducimus. \\
(1) ダミーのひとつめ。 \\
(2) ダミーのふたつめ。

\begin{thebibliography}{99}

\bibitem{Dummy-1}
名無しの権兵衛, \JPSJ {\bf 99} (2020) 012345.

\bibitem{Dummy-2}
名無野 三郎 and 何野 某, \PTP {\bf 100} (2021) 067890;\\
同上, {\it ダミー論文集 }, (架空出版, 2022) p.~1.

\bibitem{Dummy-3}
何野 某, ``ダミーの表題'', arXiv:0000.00000 (2023).

\end{thebibliography}

%%%%%%%%%%%%%%%
%%%%%%%%%%%%

\end{document}
